% !Mode:: "TeX:UTF-8"
\chapter*{致谢}
\addcontentsline{toc}{chapter}{致谢}
论文写到这里，忽然有些恍惚。光标在“致谢”两个字后面一闪一闪，像在催促，又像在等待。我坐了很久，不知从何说起。四年的时光压缩成薄薄一册，而真正厚重的部分，从来不在纸上。

首先，要感谢我的导师孟德元老师。从选题时的茫然，到中途的技术细节推进，再到最后逐字逐句的打磨，孟老师始终以极大的耐心对我进行无私的指导与帮助。您从不急于给我答案，而是用提问引我往前走。那些看似平淡的对话里，藏着您对学术的敬畏和对学生的体恤。孟老师对待工作极其严谨，对时间的珍视几乎到了苛刻的地步，但这让我看见一个认真的学者该有的样子，也让我相信，珍惜时间，扎实一点，这条路是可以走通的。未来博士生涯，我将怀着更高的热情加倍努力，追随榜样的步伐，不断前行。

其次，感谢我的指导教师冯伟老师。毕设的整个过程，冯老师始终关心着我的进度和状态，在我松懈时及时提醒，在我卡壳时耐心询问。那些定期的询问和叮嘱，看似平常，回过头看，每一句都是一种踏实的支撑。感谢师兄王琛超、王雯显，在毕设期间给予我的热心帮助。许多具体的问题，当我自己琢磨不清时，总是耐心地指导我，放下手头的事情帮我分析、给我演示。还要感谢在大学期间所有给予我教诲的老师们。课堂的知识为我打开了专业的大门，也让我在困惑时得以窥见更开阔的风景。

感谢一路同行的伙伴们。感谢一起学习的同学们，感谢生活中相互照顾的舍友们，感谢兴趣相投一起运动的朋友们。点点滴滴，都是我大学生活里最珍贵的部分。

最后，感谢我的家人。来自远方的关心从来润物无声，这份无条件的托底，是我所有底气的来源。

一段路走到这里，不是结束，是换了一种方式继续。我把所有感谢收进行囊，轻轻合上这四年。从此山高水长，愿我们都在各自的路上，走得踏实而明亮。
\cleardoublepage
